% !TEX root = main.tex

\section{Background and Related Work}
\label{sec:background}

\subsection{Problem Formulation and Physics-Informed Loss}
A general non-linear Volterra Integro-Differential Equation (VIDE) defined on the standard bounded domain $\Omega = [a, b]$ is formulated as:
\begin{equation}
    \begin{aligned}
        \mathcal{D}(u)(x) + \mathcal{F}(u)(x) 
        = \mathcal{S}(x) + \int_{a}^{x} K(x, y) \, \zeta\big(u(y)\big) \, dy,
    \end{aligned}
    \label{eq:general_vide}
\end{equation}
where $\mathcal{D}(u)(x)$ represents a linear differential operator, $\mathcal{F}(u)(x)$ denotes local algebraic state interactions, $K(x, y)$ is a continuous memory kernel, $\zeta(\cdot)$ is a non-linear mapping, and $\mathcal{S}(x)$ is a known driving source. 

To train a neural approximator $\hat{u}_\theta(x)$, we minimize the Integro-Differential Equation (IDE) residual $ \mathcal{R}(x)$ of Eq.~\eqref{eq:general_vide}. Evaluating this residual requires computing the Volterra integral, which we discretize via an $N_q$-point Gauss-Legendre quadrature scheme~\cite{moghaddam2025risn,aghaei2024pinnies}, yielding the computable residual $\mathcal{R}(x)$:
\begin{equation}
    \begin{aligned}
        \mathcal{R}(x) = \;&\mathcal{D}(\hat{u}_\theta)(x) + \mathcal{F}(\hat{u}_\theta)(x) \\
        &- \sum_{j=1}^{N_q} w_j(x) K(x, y_j) \, \zeta\big(\hat{u}_\theta(y_j)\big) - \mathcal{S}(x),
    \end{aligned}
    \label{eq:coupled_residual}
\end{equation}
where $y_j$ and $w_j(x)$ represent the quadrature nodes and integration weights. The network parameters $\theta$ are optimized by minimizing the Mean Squared Residual $\mathcal{L}_{\text{IDE}} = \frac{1}{N_r} \sum_{i=1}^{N_r} | \mathcal{R}(x_i) |^2$ over a batch of $N_r$ collocation points.

\subsection{Volterra Solvers and Manufactured Solutions}

In recent years, several specialized deep learning architectures have been introduced to solve integral and integro-differential equations, including Auxiliary PINNs (A-PINN)~\cite{yuan2022apinn}, and internal variable reformulations for exponential kernels~\cite{Jang_2026}. Notably, Moghaddam et al.~\cite{moghaddam2025risn} proposed the Residual Integration Solver Network (RISN) to mitigate gradient degradation in integral operators.

Unlike standard machine learning tasks, forward operator learning often lacks experimental observational datasets. To rigorously validate architectures without real-world data, we employ the Method of Manufactured Solutions (MMS)~\cite{moghaddam2025risn}. In MMS, a smooth analytical target solution $u_{\text{exact}}(x)$ is selected \textit{a priori} and substituted into the integro-differential operator to analytically synthesize the corresponding source term $\mathcal{S}(x)$. 
This guarantees to have an exact objective to use for precise and reliable error quantification against baseline architectures.

While RISN and related variants provide structural improvements over standard baseline Multi-Layer Perceptrons, applying an MLP-based architecture to solve Eq.~\eqref{eq:coupled_residual} introduces a structural mismatch, which we refer to as the \textit{pointwise paradox}. Volterra equations are intrinsically non-local, yet MLPs evaluate each coordinate in strict isolation, delegating the entire sequence-awareness to the external loss function.

We hypothesize that this structural mismatch gives rise to an important optimization pathology, which we call the \textit{gradient conflict}. As shown in Eq.~\eqref{eq:coupled_residual}, computing the numerical residual $\mathcal{R}(x)$ mathematically couples the network's prediction at the present coordinate, $\hat{u}_\theta(x)$, with its predictions at historical quadrature nodes, $\hat{u}_\theta(y_j)$. During backpropagation, the optimizer must update the shared global weights $\theta$ based on these coupled outputs. Consequently, the gradients computed for the local state $\hat{u}_\theta(x)$ can directly conflict with those required to minimize approximation errors across past states $\hat{u}_\theta(y_j)$, trapping the optimizer in suboptimal local minima.

\subsection{Resolving the Paradox via Attention}
Tom mitigate this potential gradient conflict, we propose embedding a Cross-Attention mechanism into the physics-informed architecture. Because cross-attention is intrinsically set-aware, it allows the model to correlate a specific query coordinate (the present state) with a structured latent basis (the historical memory) during the forward pass, embedding memory integration directly into the network architecture.